\cvsection{Open-Source and Artifacts}
\cvitem{\textbf{Agentic cryptographic programming:} \href{https://fhe-coder.github.io/}{FHE-Coder benchmark and project resources}.}
\cvitem{\textbf{Encrypted arithmetic and comparison:} \href{https://github.com/UCF-Lou-Lab-PET/Universal-BGV}{Universal-BGV}.}
\cvitem{\textbf{Private verification:} \href{https://github.com/UCF-Lou-Lab-PET/zkformer}{zkVC / zkformer}.}
\cvitem{\textbf{Privacy-preserving computing and AI:} \href{https://github.com/UCF-Lou-Lab-PET}{UCF-Lou-Lab-PET}.}
\cvsection{Selected Research Contributions and Impact}
\begin{itemize}
\item \textbf{Private computation and verification.} SHE (NeurIPS 2019) reduces encrypted neural-network computation costs. Universal-BGV (IEEE S\&P 2026) connects arithmetic and comparisons within one encryption scheme. DataSeal (IEEE S\&P 2025) checks encrypted computation results, while zkVC (DAC 2025) accelerates private verification through zero-knowledge proofs.
\item \textbf{Compilers for secure AI.} CipherFlow (NeurIPS 2026) coordinates secure protocols and hardware choices to reduce inference latency. HEMET (ICML 2021) and SAFENet (ICLR 2021) adapt AI models to the costs of secure computation.
\item \textbf{Protected learning and applications.} DictPFL (NeurIPS 2025) supports federated learning with encrypted gradients. SecureRoute (DAC 2026) protects model selection and inference. FHAIM (ICML 2026) combines encrypted computation with differential privacy for synthetic data generation.
\item \textbf{AI for cryptographic programming.} FHE-Coder (ICLR 2026) helps generate cryptographic programs using expert knowledge and external checks of functionality and security, reducing the programming expertise needed to work with FHE libraries.
\end{itemize}
\cvitem{\textbf{Ongoing directions:} native-complexity, general-purpose, verifiable FHE (nuvFHE); sub-linear proofs for general matrix multiplication (redZKP); and AI agents for cryptanalysis and cryptographic optimization.}
